\documentclass[10pt]{extarticle}
\usepackage{fontspec}
\setmainfont{Leelawadee UI}
\XeTeXlinebreaklocale "th_TH"
\XeTeXlinebreakskip = 0pt plus 1pt minus 0.1pt
\usepackage[a4paper,top=1.3cm,bottom=1.5cm,left=1.4cm,right=1.4cm]{geometry}
\usepackage{amsmath,amssymb}
\usepackage{xcolor}
\usepackage{colortbl}
\usepackage{booktabs}
\usepackage{tikz}
\usepackage{pgfplots}
\pgfplotsset{compat=1.18}
\usepackage{enumitem}
\setlist{nosep,leftmargin=1.3em,itemsep=1pt,topsep=2pt}
\usepackage{titlesec}
\usepackage{tcolorbox}
\tcbuselibrary{skins,breakable}
\usepackage{array}
\usepackage{longtable}
\usepackage{multirow}

\definecolor{navy}{RGB}{20,40,90}
\definecolor{gold}{RGB}{190,140,20}
\definecolor{lightbg}{RGB}{240,244,250}
\definecolor{solbg}{RGB}{247,247,240}
\definecolor{ideabg}{RGB}{236,246,238}
\definecolor{ideafr}{RGB}{60,120,70}
\definecolor{compbg}{RGB}{236,247,247}
\definecolor{compfr}{RGB}{20,120,120}
\definecolor{qbg}{RGB}{245,240,250}
\definecolor{qfr}{RGB}{90,60,140}

\titleformat{\section}{\normalfont\bfseries\Large\color{navy}}{\thesection}{0.5em}{}[{\color{gold}\titlerule[1.2pt]}]
\titlespacing*{\section}{0pt}{10pt}{5pt}
\titleformat{\subsection}{\normalfont\bfseries\large\color{gold!70!black}}{\thesubsection}{0.4em}{}
\titlespacing*{\subsection}{0pt}{7pt}{3pt}
\titleformat{\subsubsection}{\normalfont\bfseries\color{navy}}{\thesubsubsection}{0.4em}{}
\titlespacing*{\subsubsection}{0pt}{5pt}{2pt}

\newtcolorbox{formulabox}{
  colback=lightbg, colframe=navy!60, boxrule=0.4pt, breakable,
  arc=1pt, left=4pt, right=4pt, top=3pt, bottom=3pt,
}
\newtcolorbox{solbox}{
  colback=solbg, colframe=gold!70!black, boxrule=0.4pt, breakable,
  arc=1pt, left=4pt, right=4pt, top=3pt, bottom=3pt,
}
\newtcolorbox{ideabox}{
  colback=ideabg, colframe=ideafr, boxrule=0.4pt, breakable,
  arc=1pt, left=4pt, right=4pt, top=3pt, bottom=3pt,
}
\newtcolorbox{outbox}{
  colback=compbg, colframe=compfr, boxrule=0.4pt, breakable,
  arc=1pt, left=4pt, right=4pt, top=3pt, bottom=3pt,
}
\newtcolorbox{qbox}{
  colback=qbg, colframe=qfr, boxrule=0.5pt, breakable,
  arc=1pt, left=4pt, right=4pt, top=3pt, bottom=3pt,
  fonttitle=\bfseries\color{white}, colbacktitle=qfr, title={โจทย์ (ต้นฉบับ)}
}
\newtcolorbox{stepbox}[1]{
  colback=white, colframe=navy, boxrule=0.6pt, breakable,
  arc=1pt, left=5pt, right=5pt, top=3pt, bottom=4pt,
  fonttitle=\bfseries\color{white}, colbacktitle=navy, title={#1}
}

\newcommand{\dist}[1]{\subsection*{\textcolor{navy}{\textbullet\ #1}}}
\newcommand{\case}[1]{\textbf{\textcolor{gold!70!black}{#1}}}
\newcommand{\concept}[1]{\par\vspace{2pt}\noindent{\footnotesize\textbf{\textcolor{ideafr}{\ding{228}\ แนวคิด:}} #1}\vspace{2pt}\par}
\newcommand{\ding}[1]{$\blacktriangleright$}
\newcommand{\spssout}[1]{\vspace{2pt}\noindent{\small\textbf{\textcolor{compfr}{Output -- SPSS: #1}}}\par\vspace{1pt}}
\newcommand{\excelout}[1]{\vspace{2pt}\noindent{\small\textbf{\textcolor{compfr}{Output -- Excel: #1}}}\par\vspace{1pt}}

\setlength{\parindent}{0pt}
\setlength{\parskip}{1pt}
\renewcommand{\arraystretch}{1.2}

\begin{document}
\begin{center}
{\LARGE\bfseries\color{navy} Chapter 5 (Part 2): การวิเคราะห์ความแปรปรวน (ANOVA) และการทดสอบไคกำลังสอง (Chi-Square Test)}\\[3pt]
{\normalsize สรุปเนื้อหาพร้อมตัวอย่างและ Output SPSS/Excel ฉบับเต็ม — อ.ดร.ธนวรรณ ประฮาดไชย \quad Department of Statistics, Faculty of Science, KKU}
\end{center}
\vspace{3pt}
{\color{gold}\hrule height 1.6pt}
\vspace{6pt}

\section{ภาพรวม: ANOVA vs Chi-square Test}

ทั้งสองหัวข้อในบทนี้เป็นการทดสอบสมมติฐานที่ขยายจาก Part 1 ไปสู่กรณีที่มี \textbf{มากกว่า 2 กลุ่ม} หรือ \textbf{ตัวแปรเชิงคุณภาพ 2 ตัวแปร} โดยมีความแตกต่างกันดังนี้

\begin{center}\small
\begin{tabular}{@{}lll@{}}
\toprule
\rowcolor{lightbg}\textbf{ประเด็น} & \textbf{One-way ANOVA} & \textbf{Chi-square test}\\
\midrule
ตัวแปรตาม & เชิงปริมาณ & เชิงจำแนก\\
ตัวแปรอิสระ & เชิงจำแนก & เชิงจำแนก\\
สิ่งที่เปรียบเทียบ & ค่าเฉลี่ย & ความถี่ / สัดส่วน\\
สถิติทดสอบ & $F$ & $\chi^2$\\
คำถามหลัก & ค่าเฉลี่ยต่างกันไหม & ตัวแปรสัมพันธ์กันไหม\\
ตัวอย่างผลลัพธ์ & $\mu_1\ne\mu_2$ & ไม่เป็นอิสระ\\
\bottomrule
\end{tabular}
\end{center}
\concept{หลักการเลือกว่าจะใช้ ANOVA หรือ Chi-square ให้ดูที่ \textbf{มาตราวัดของตัวแปรตาม} เป็นหลัก: ถ้าตัวแปรตามเป็นตัวเลขต่อเนื่อง (เช่น คะแนน, น้ำหนัก) และต้องการเปรียบเทียบค่าเฉลี่ยระหว่าง $\ge3$ กลุ่ม $\to$ ใช้ ANOVA; ถ้าตัวแปรตามเป็นข้อมูลเชิงกลุ่ม/นับความถี่ (เช่น ระดับ, ประเภท) และต้องการดูความสัมพันธ์หรือความแตกต่างของสัดส่วน $\to$ ใช้ Chi-square}

\clearpage
\section{การวิเคราะห์ความแปรปรวนแบบจำแนกทางเดียว (One-Way ANOVA)}

\textbf{แนวคิด:} ใช้ทดสอบว่าค่าเฉลี่ยของ\textbf{ประชากรตั้งแต่ 2 กลุ่มขึ้นไป}ที่เป็นอิสระต่อกัน (สุ่มมาจากประชากรคนละกลุ่ม) แตกต่างกันหรือไม่ เป็นการขยายแนวคิดจากการทดสอบผลต่างค่าเฉลี่ย 2 ประชากรใน Part 1 ไปสู่กรณีที่มีมากกว่า 2 กลุ่ม (ถ้ามีแค่ 2 กลุ่ม การใช้ ANOVA จะให้ผลตรงกับการทดสอบ $T$ อิสระแบบสองทางทุกประการ)
\concept{เหตุผลที่ต้อง ``ANOVA'' (วิเคราะห์ความแปรปรวน) ทั้ง ๆ ที่คำถามคือเปรียบเทียบ \textbf{ค่าเฉลี่ย}: เพราะถ้าเปรียบเทียบค่าเฉลี่ยทีละคู่ด้วย $T$-test หลาย ๆ ครั้ง จะทำให้ความคลาดเคลื่อนชนิดที่ 1 สะสมเกินกว่า $\alpha$ ที่ตั้งไว้ (Type I error inflation) ANOVA จึงใช้วิธีเปรียบเทียบ \textbf{ความแปรปรวนระหว่างกลุ่ม} กับ \textbf{ความแปรปรวนภายในกลุ่ม} พร้อมกันในการทดสอบครั้งเดียว ถ้าความแปรปรวนระหว่างกลุ่มมากกว่าความแปรปรวนภายในกลุ่มอย่างมีนัยสำคัญ แสดงว่าอย่างน้อยมีค่าเฉลี่ยบางกลุ่มที่ต่างจากกลุ่มอื่น}

\dist{ข้อตกลงเบื้องต้น (Assumptions) -- ต้องตรวจก่อนทุกครั้ง}
\begin{enumerate}
\item ตัวอย่างทุกกลุ่มถูกสุ่มมาจากประชากรที่\textbf{มีการแจกแจงปกติ} -- ทดสอบด้วย Kolmogorov--Smirnov ($n>50$) หรือ Shapiro--Wilk ($n\le50$)
$$H_0:\text{ข้อมูลมีการแจกแจงแบบปกติ}\qquad H_1:\text{ข้อมูลไม่มีการแจกแจงแบบปกติ}$$
\item ตัวอย่างทุกกลุ่มถูกสุ่มมาจากประชากรที่มี\textbf{ความแปรปรวนไม่ต่างกัน} -- ทดสอบด้วย Levene's Test
$$H_0:\text{ประชากรทุกกลุ่มมีความแปรปรวนเท่ากัน}\qquad H_1:\text{ประชากรอย่างน้อย 1 กลุ่มมีความแปรปรวนต่างจากกลุ่มอื่น}$$
\end{enumerate}
คำสั่งใน SPSS: \texttt{Analyze > Descriptive Statistics > Explore...} (เลือกตัวแปรตามใส่ Dependent List และตัวแปรจำแนกกลุ่มใส่ Factor List)

\dist{การทดสอบสมมติฐาน}
$$H_0:\mu_1=\mu_2=\cdots=\mu_k \qquad H_1:\mu_i\text{ อย่างน้อย 2 ค่าไม่เท่ากัน},\ i=1,2,\dots,k$$
กำหนดระดับนัยสำคัญ $(\alpha)$ แล้วคำนวณตัวสถิติทดสอบจากตาราง ANOVA ดังนี้

\begin{center}\small
\begin{tabular}{@{}lccccc@{}}
\toprule
\rowcolor{lightbg}\textbf{แหล่งความผันแปร (SOV)} & \textbf{df} & \textbf{SS} & \textbf{MS} & \textbf{F}\\
\midrule
สิ่งทดลอง (Treatment) & $k-1$ & $SSTr$ & $MSTr=\dfrac{SSTr}{k-1}$ & \multirow{2}{*}{$F=\dfrac{MSTr}{MSE}$}\\
ความคลาดเคลื่อน (Error) & $n-k$ & $SSE$ & $MSE=\dfrac{SSE}{n-k}$ &\\
รวม (Total) & $n-1$ & $SST$ & & \\
\bottomrule
\end{tabular}
\end{center}

\begin{formulabox}
\textbf{สูตรผลรวมกำลังสอง} ($X_{ij}=$ ค่าสังเกตที่ $j$ ของสิ่งทดลองที่ $i$, $X_{i.}=$ ผลรวมของสิ่งทดลองที่ $i$, $X_{..}=$ ผลรวมทั้งหมด, $n=\sum n_i$)
$$SST=\sum_{i=1}^k\sum_{j=1}^{n_i}X_{ij}^2-\frac{X_{..}^2}{n} \qquad
SSTr=\sum_{i=1}^k\frac{X_{i.}^2}{n_i}-\frac{X_{..}^2}{n} \qquad
SSE=SST-SSTr$$
\end{formulabox}
\concept{เกณฑ์ตัดสินใจเหมือนบทที่แล้วทุกประการ: ปฏิเสธ $H_0$ ถ้า p-value $\le\alpha$ (เทียบเท่ากับ $F$ คำนวณได้ $>F$ วิกฤต) เมื่อปฏิเสธ $H_0$ สรุปได้เพียงว่า ``มีอย่างน้อย 2 กลุ่มที่ค่าเฉลี่ยต่างกัน'' เท่านั้น \textbf{ยังไม่รู้ว่าคู่ไหนต่างกันบ้าง} จึงต้องทำการเปรียบเทียบเชิงพหุคูณ (Multiple Comparison) ต่อเป็นขั้นที่ 2 เสมอ}

\dist{การเปรียบเทียบเชิงพหุคูณ (Multiple Comparison / Post Hoc Test)}
ทำต่อเมื่อผล ANOVA \textbf{ปฏิเสธ} $H_0$ เท่านั้น เพื่อหาว่าค่าเฉลี่ยคู่ใดบ้างที่แตกต่างกันจริง
\begin{center}\small
\begin{tabular}{@{}p{2.6cm}p{9.5cm}@{}}
\toprule
\rowcolor{lightbg}\textbf{วิธี} & \textbf{เกณฑ์การเลือกใช้}\\
\midrule
Least Significant Difference (LSD) & เปรียบเทียบค่าเฉลี่ยประชากรได้ครั้งละหลายคู่ (เสรีสุด แต่ Type I error สะสมได้ง่ายที่สุด)\\
Tukey's Honestly Significant Difference (HSD) & ใช้เมื่อตัวอย่างแต่ละกลุ่มมีขนาด\textbf{เท่ากัน}\\
Scheffe's Post Hoc Comparison & ใช้ได้ทั้งกรณีตัวอย่างแต่ละกลุ่มมีขนาด\textbf{เท่ากันหรือไม่เท่ากัน}ก็ได้ (เข้มงวดที่สุด)\\
Bonferroni Test & ใช้เมื่อจำนวนกลุ่มตัวอย่าง\textbf{ไม่เกิน 5 กลุ่ม}\\
\bottomrule
\end{tabular}
\end{center}
\concept{ทั้ง 4 วิธีนี้ล้วนแก้ปัญหาเดียวกันคือการควบคุม Type I error ที่สะสมจากการเปรียบเทียบหลายคู่พร้อมกัน แต่ต่างกันที่ระดับความเข้มงวด/เงื่อนไขการใช้งาน โจทย์ในเอกสารนี้เลือกใช้ \textbf{Scheffe} (กรณีตัวอย่างขนาดไม่เท่ากันในกรณีศึกษาความเครียด) และ \textbf{LSD} (กรณีตัวอย่างขนาดเท่ากันในตัวอย่างอาหารสุกร) เป็นตัวอย่าง}

\clearpage
\section{ตัวอย่างที่ 1 (ANOVA): กรณีศึกษาความเครียดของนักศึกษาจากโควิด-19}

\begin{qbox}
\textbf{งานวิจัยเรื่อง:} ผลกระทบจากการแพร่ระบาดของโรคโควิด-19 ต่อความเครียดของนักศึกษาระดับปริญญาตรี ($n=150$ คน)\\
อยากทราบว่า นักศึกษาแต่ละชั้นปี (ปี 1--4) มีคะแนนความเครียดเฉลี่ย แตกต่างกันหรือไม่?? ที่ระดับนัยสำคัญ 0.05
\end{qbox}

\begin{stepbox}{ขั้นที่ 1: กำหนดสมมติฐาน}
$$H_0:\mu_1=\mu_2=\mu_3=\mu_4 \qquad H_1:\mu_i\text{ อย่างน้อย 2 ค่าไม่เท่ากัน},\ i=1,2,3,4$$
\concept{เปรียบเทียบค่าเฉลี่ยคะแนนความเครียดระหว่าง \textbf{4 กลุ่ม} (ชั้นปี 1--4) ที่เป็นอิสระต่อกัน (นักศึกษาคนละคนกัน) จึงเป็นกรณี One-Way ANOVA ไม่ใช่ $T$-test อีกต่อไปเพราะมีมากกว่า 2 กลุ่ม}
\end{stepbox}

\begin{stepbox}{ขั้นที่ 2: กำหนดระดับนัยสำคัญ}
กำหนดระดับนัยสำคัญที่ $\alpha=0.05$
\end{stepbox}

\begin{stepbox}{ขั้นที่ 3: พิจารณาข้อตกลงเบื้องต้น}
ขนาดตัวอย่างแต่ละชั้นปี: ปี 1 $=100$, ปี 2 $=15$, ปี 3 $=11$, ปี 4 $=24$ ($n=150$)

\spssout{Tests of Normality \& Test of Homogeneity of Variance}
\begin{center}\scriptsize
\begin{tabular}{@{}lcccccc@{}}
\toprule
& \multicolumn{3}{c}{Kolmogorov-Smirnov} & \multicolumn{3}{c}{Shapiro-Wilk}\\
\cmidrule(lr){2-4}\cmidrule(lr){5-7}
ชั้นปีที่เรียน & Statistic & df & Sig. & Statistic & df & Sig.\\
\midrule
1 & .058 & 100 & .200* & .988 & 100 & .503\\
2 & .205 & 15 & .092 & .928 & 15 & .253\\
3 & .165 & 11 & .200* & .919 & 11 & .311\\
4 & .192 & 24 & .023 & .919 & 24 & .056\\
\bottomrule
\end{tabular}
\end{center}
{\scriptsize Levene Statistic (Based on Mean) $=1.753$, $df1=3$, $df2=146$, Sig.$=.159$}

พบว่าทุกชั้นปีมี p-value (Shapiro--Wilk) $>0.05\Rightarrow$ \textbf{ยอมรับ} $H_0$ ว่าข้อมูลมีการแจกแจงปกติ และ Levene's Sig.$=.159>0.05\Rightarrow$ \textbf{ยอมรับ} $H_0$ ว่าความแปรปรวนของทั้ง 4 กลุ่มไม่ต่างกัน $\Rightarrow$ เข้าเงื่อนไขครบทั้ง 2 ข้อ จึงใช้ One-Way ANOVA ได้
\concept{ชั้นปี 4 มี K--S Sig.$=.023<0.05$ (ดูเหมือนผิดปกติ) แต่ S--W Sig.$=.056>0.05$ (ยอมรับปกติ) -- เมื่อผลทั้งสองสถิติขัดแย้งกันเล็กน้อยแบบนี้ ให้ยึดตามหลักที่กำหนดไว้ตั้งแต่ต้น คือ ใช้ Kolmogorov--Smirnov เมื่อ $n>50$ และ Shapiro--Wilk เมื่อ $n\le50$ -- เนื่องจากชั้นปี 4 มี $n=24\le50$ จึงใช้ผลจาก Shapiro--Wilk ($.056>0.05$) เป็นเกณฑ์ตัดสิน ไม่ใช่ K--S}
\end{stepbox}

\begin{stepbox}{ขั้นที่ 4: คำนวณค่าสถิติและหาค่า p-value}
คำสั่งใน SPSS: \texttt{Analyze > Compare Means > One-Way ANOVA...}

\spssout{Descriptives}
\begin{center}\scriptsize
\begin{tabular}{@{}lcccccc@{}}
\toprule
ชั้นปีที่เรียน & N & Mean & Std. Deviation & Std. Error & Minimum & Maximum\\
\midrule
1 & 100 & 20.93 & 6.921 & .692 & 6 & 36\\
2 & 15 & 20.13 & 6.468 & 1.670 & 4 & 32\\
3 & 11 & 29.45 & 5.373 & 1.620 & 22 & 37\\
4 & 24 & 27.79 & 4.995 & 1.020 & 21 & 40\\
Total & 150 & 22.57 & 7.200 & .588 & 4 & 40\\
\bottomrule
\end{tabular}
\end{center}

\spssout{ANOVA}
\begin{center}\small
\begin{tabular}{@{}lccccc@{}}
\toprule
& Sum of Squares & df & Mean Square & F & Sig.\\
\midrule
Between Groups & 1533.764 & 3 & 511.255 & 12.057 & $<.001$\\
Within Groups & 6190.929 & 146 & 42.404 & & \\
Total & 7724.693 & 149 & & & \\
\bottomrule
\end{tabular}
\end{center}

\excelout{ยืนยันผลด้วย Data Analysis: Anova Single Factor}
\begin{center}\scriptsize
\begin{tabular}{@{}lcccccc@{}}
\toprule
Groups & Count & Sum & Average & Variance & & \\
\midrule
ปี 1 & 100 & 2093 & 20.93 & 47.904 & & \\
ปี 2 & 15 & 302 & 20.133 & 41.838 & & \\
ปี 3 & 11 & 324 & 29.455 & 28.873 & & \\
ปี 4 & 24 & 667 & 27.792 & 24.955 & & \\
\midrule
Source & SS & df & MS & F & P-value & F crit\\
Between Groups & 1533.764 & 3 & 511.255 & 12.057 & 0.0000004 & 2.667\\
Within Groups & 6190.929 & 146 & 42.404 & & & \\
\bottomrule
\end{tabular}
\end{center}
\concept{Excel และ SPSS ให้ผลตรงกันทุกตัวเลข (ต่างเครื่องมือ แต่คำนวณจากสูตรเดียวกัน) -- นี่เป็นวิธีที่ดีในการตรวจทานความถูกต้องของการคำนวณ Excel รายงาน P-value แบบเต็ม ($0.0000004$) ในขณะที่ SPSS ตัดทอนเป็น ``$<.001$'' ซึ่งสื่อความหมายเดียวกัน}
\end{stepbox}

\begin{stepbox}{ขั้นที่ 5: ทำการตัดสินใจ และสรุปผลการทดสอบ}
p-value $<.001<\alpha=0.05\Rightarrow$ \textbf{ปฏิเสธ} $H_0$

สรุป: นักศึกษาแต่ละชั้นปี มีคะแนนความเครียดเฉลี่ย\textbf{แตกต่างกัน} ที่ระดับนัยสำคัญ 0.05
\concept{เมื่อปฏิเสธ $H_0$ ของ ANOVA แล้ว จะรู้แค่ว่า ``มีอย่างน้อย 2 ชั้นปีที่ค่าเฉลี่ยต่างกัน'' แต่\textbf{ยังบอกไม่ได้ว่าคู่ไหนบ้าง} จึงต้องทำ Multiple Comparison ต่อทันที (ขั้นตอนถัดไป) ก่อนจะสรุปผลอย่างสมบูรณ์}
\end{stepbox}

\dist{การเปรียบเทียบพหุคูณ (Scheffe's Method) เพื่อหาว่าชั้นปีใดต่างกัน}
\spssout{Post Hoc Tests -- Multiple Comparisons (Scheffe), Dependent Variable: คะแนนความเครียด}
\begin{center}\scriptsize
\begin{tabular}{@{}llccc@{}}
\toprule
(I) ชั้นปี & (J) ชั้นปี & Mean Difference (I--J) & Std. Error & Sig.\\
\midrule
1 & 2 & .797 & 1.803 & .978\\
1 & 3 & $-8.525^*$ & 2.069 & .001\\
1 & 4 & $-6.862^*$ & 1.480 & $<.001$\\
2 & 3 & $-9.321^*$ & 2.585 & .006\\
2 & 4 & $-7.658^*$ & 2.143 & .006\\
3 & 4 & 1.663 & 2.371 & .920\\
\bottomrule
\end{tabular}
\end{center}

\spssout{Homogeneous Subsets (Scheffe), $\alpha=0.05$}
\begin{center}\small
\begin{tabular}{@{}lcccc@{}}
\toprule
ชั้นปีที่เรียน & N & Subset 1 & Subset 2\\
\midrule
2 & 15 & 20.13 & \\
1 & 100 & 20.93 & \\
4 & 24 & & 27.79\\
3 & 11 & & 29.45\\
\midrule
Sig. & & .986 & .891\\
\bottomrule
\end{tabular}
\end{center}

จากผลเปรียบเทียบคู่ (I)--(J) พบว่า ปี 1 กับ ปี 2 มี Sig.$=.978>0.05\Rightarrow$ ค่าเฉลี่ยความเครียด\textbf{ไม่แตกต่างกัน}; ปี 3 กับ ปี 4 มี Sig.$=.920>0.05\Rightarrow$ ค่าเฉลี่ย\textbf{ไม่แตกต่างกัน}เช่นกัน; แต่ทุกคู่ระหว่างกลุ่ม \{ปี 1, ปี 2\} กับกลุ่ม \{ปี 3, ปี 4\} มี Sig.$\le0.05\Rightarrow$ \textbf{แตกต่างกันอย่างมีนัยสำคัญ}ทั้งหมด

สรุป: นักศึกษาแบ่งได้เป็น 2 กลุ่มค่าเฉลี่ย คือ กลุ่มความเครียดต่ำ (ปี 1 $=20.93$, ปี 2 $=20.13$) และกลุ่มความเครียดสูง (ปี 3 $=29.45$, ปี 4 $=27.79$) โดยรุ่นพี่ปี 3--4 มีความเครียดเฉลี่ยสูงกว่ารุ่นน้อง ปี 1--2 อย่างมีนัยสำคัญ
\concept{ตาราง Homogeneous Subsets เป็นวิธีอ่านผล Multiple Comparison ที่รวดเร็วกว่าการไล่อ่านทีละคู่: ชั้นปีที่อยู่ใน\textbf{คอลัมน์ (Subset) เดียวกัน}คือกลุ่มที่ค่าเฉลี่ย\textbf{ไม่แตกต่างกัน} ส่วนชั้นปีที่อยู่คนละ Subset คือกลุ่มที่ค่าเฉลี่ย\textbf{แตกต่างกัน} -- ในที่นี้ Subset 1 = \{ปี 2, ปี 1\} และ Subset 2 = \{ปี 4, ปี 3\} ไม่มีชั้นปีใดอยู่ซ้อนทั้ง 2 Subset แสดงว่าผลแบ่งกลุ่มชัดเจนมาก}

\clearpage
\section{ตัวอย่างที่ 2 (ANOVA): ประสิทธิภาพอาหารสุกร 4 สูตร}

\begin{qbox}
ในการศึกษาประสิทธิภาพอาหารที่ใช้เลี้ยงสุกร 4 สูตร โดยแต่ละสูตรใช้เลี้ยงลูกสุกรสายพันธุ์เดียวกันที่มีอายุและน้ำหนักเท่ากัน โดยสุ่มรวม 5 ตัวต่อสูตร หลังการทดลองชั่งน้ำหนัก (กก.) ปรากฏผลดังนี้ จงทดสอบว่าอาหารสุกรทั้ง 4 สูตร มีประสิทธิภาพแตกต่างกันหรือไม่ ที่ระดับนัยสำคัญ 0.05

\begin{center}\small
\begin{tabular}{@{}c cccc@{}}
\toprule
& \multicolumn{4}{c}{สูตรอาหารที่}\\
\cmidrule(lr){2-5}
& 1 & 2 & 3 & 4\\
\midrule
& 137.8 & 154.4 & 163.8 & 143.2\\
& 129.3 & 146.3 & 165.6 & 141.0\\
& 146.1 & 143.8 & 162.4 & 150.7\\
& 132.9 & 149.5 & 157.3 & 137.8\\
& 139.7 & 159.6 & 161.0 & 145.5\\
\midrule
รวม ($X_{i.}$) & 685.8 & 753.6 & 810.1 & 718.2\\
\bottomrule
\end{tabular}\\[2pt]
\footnotesize รวมทั้งหมด $X_{..}=2{,}967.7$, $n=20$
\end{center}
\end{qbox}

\begin{stepbox}{ขั้นที่ 1--3: สมมติฐาน ระดับนัยสำคัญ และตัวสถิติทดสอบ}
$$H_0:\mu_1=\mu_2=\mu_3=\mu_4 \qquad H_1:\mu_i\text{ อย่างน้อย 2 ค่าไม่เท่ากัน},\ i=1,2,\dots,4$$
กำหนดระดับนัยสำคัญ $\alpha=0.05$; ใช้ตัวสถิติ $F=MSTr/MSE$ จากตาราง ANOVA
\concept{โจทย์นี้ตัวอย่างแต่ละกลุ่มมีขนาด\textbf{เท่ากัน} ($n_i=5$ ทุกสูตร) ต่างจากตัวอย่างที่ 1 (ขนาดกลุ่มไม่เท่ากัน) จุดสังเกตนี้จะมีผลตอนเลือกวิธี Multiple Comparison ในขั้นถัดไป (จะเลือกใช้ LSD ได้ง่ายเพราะสูตร LSD คำนวณตรงไปตรงมาเมื่อขนาดกลุ่มเท่ากัน)}
\end{stepbox}

\begin{stepbox}{ขั้นที่ 4: คำนวณค่าสถิติด้วยมือ}
$$SST=\sum\sum X_{ij}^2-\frac{X_{..}^2}{n}=(137.8)^2+(129.3)^2+\cdots+(145.5)^2-\frac{(2{,}967.7)^2}{20}=2{,}161.45$$
$$SSTr=\sum_{i=1}^4\frac{X_{i.}^2}{n_i}-\frac{X_{..}^2}{n}=\frac{(685.8)^2}{5}+\frac{(753.6)^2}{5}+\frac{(810.1)^2}{5}+\frac{(718.2)^2}{5}-\frac{(2{,}967.7)^2}{20}=1{,}699.41$$
$$SSE=SST-SSTr=2{,}161.45-1{,}699.41=462.04$$

\begin{center}\small
\begin{tabular}{@{}lcccc@{}}
\toprule
แหล่งความแปรปรวน & df & SS & MS & F\\
\midrule
สูตรอาหาร & 3 & 1,699.41 & 566.47 & 19.61\\
ความคลาดเคลื่อน & 16 & 462.04 & 28.88 & \\
รวม & 19 & 2,161.45 & & \\
\bottomrule
\end{tabular}
\end{center}
$$F=\frac{MSTr}{MSE}=\frac{566.47}{28.88}=19.61 \qquad p\text{-value}=P(F_{(3,16)}>19.61)=0.0000312$$
\concept{$df$ ของสูตรอาหาร $=k-1=4-1=3$, $df$ ของความคลาดเคลื่อน $=n-k=20-4=16$ ตรงกับจำนวนสุกรทั้งหมดลบจำนวนสูตร นี่คือตัวอย่างการคำนวณ ANOVA ด้วยมือแบบเต็มขั้นตอน ซึ่งในทางปฏิบัติ SPSS/Excel จะคำนวณให้อัตโนมัติ (ดูตัวอย่างที่ 1 ที่ใช้ทั้งสองโปรแกรม) แต่การฝึกคำนวณด้วยมือช่วยให้เข้าใจที่มาของตัวเลขในตาราง Output}
\end{stepbox}

\begin{stepbox}{ขั้นที่ 5: ทำการตัดสินใจ และสรุปผลการทดสอบ}
p-value $=0.0000312<\alpha=0.05\Rightarrow$ \textbf{ปฏิเสธ} $H_0$

สรุป: อาหารสุกรอย่างน้อย 2 สูตร มีประสิทธิภาพ\textbf{แตกต่างกัน} ที่ระดับนัยสำคัญ 0.05
\end{stepbox}

\dist{การเปรียบเทียบพหุคูณด้วยวิธี LSD}
\begin{formulabox}
$$LSD=t_{\alpha/2,\,v}\sqrt{MSE\left(\frac1{n_i}+\frac1{n_j}\right)}, \qquad v=\sum n_i-k$$
เกณฑ์: ถ้า $|\bar X_{i.}-\bar X_{j.}|\ge LSD$ แสดงว่าคู่นั้นแตกต่างอย่างมีนัยสำคัญ; ถ้าน้อยกว่า แสดงว่าไม่แตกต่าง
\end{formulabox}

ที่ $\alpha=0.05$, $df=16$: $t_{0.025,16}=2.12$, $MSE=28.88$
$$LSD=2.12\sqrt{28.88\left(\frac15+\frac15\right)}=7.21$$

\begin{center}\small
\begin{tabular}{@{}lccc@{}}
\toprule
เปรียบเทียบสูตรอาหาร & $\bar X_{i.}-\bar X_{j.}$ & $|\bar X_{i.}-\bar X_{j.}|$ & ผลการทดสอบ\\
\midrule
1 กับ 2 & $137.16-150.72$ & $13.56>7.21$ & มีนัยสำคัญ\\
1 กับ 3 & $137.16-162.02$ & $24.86>7.21$ & มีนัยสำคัญ\\
1 กับ 4 & $137.16-143.64$ & $6.48<7.21$ & \textbf{ไม่มีนัยสำคัญ}\\
2 กับ 3 & $150.72-162.02$ & $11.30>7.21$ & มีนัยสำคัญ\\
2 กับ 4 & $150.72-143.64$ & $7.08<7.21$ & \textbf{ไม่มีนัยสำคัญ}\\
3 กับ 4 & $162.02-143.64$ & $18.38>7.21$ & มีนัยสำคัญ\\
\bottomrule
\end{tabular}
\end{center}

สรุปผลด้วยวิธีลากเส้น (ที่ $\alpha=0.05$):
$$\underbrace{\text{สูตรที่ 1}\quad\text{สูตรที่ 4}}_{\text{ไม่ต่างกัน}}\quad\underbrace{\text{สูตรที่ 4}\quad\text{สูตรที่ 2}}_{\text{ไม่ต่างกัน}}\quad\text{สูตรที่ 3}$$
\concept{ผลลัพธ์นี้อ่านได้ว่า สูตร 1 กับสูตร 4 น้ำหนักเฉลี่ยไม่ต่างกัน และสูตร 4 กับสูตร 2 ก็ไม่ต่างกัน แต่สูตร 1 กับสูตร 2 กลับ\textbf{ต่างกัน} -- ลักษณะการ ``เหลื่อมกลุ่ม'' แบบนี้ (สูตร 4 เชื่อมโยงทั้งสองกลุ่มได้) เป็นเรื่องปกติของ Multiple Comparison เพราะแต่ละคู่ทดสอบแยกกัน ไม่ได้บังคับให้แบ่งเป็นกลุ่มที่ไม่ซ้อนทับกันเหมือนตัวอย่างที่ 1 ส่วนสูตรที่ 3 (น้ำหนักเฉลี่ยสูงสุด $=162.02$) แตกต่างจากทุกสูตรอย่างมีนัยสำคัญ จึงเป็นสูตรที่มีประสิทธิภาพโดดเด่นที่สุด}

\clearpage
\section{ข้อมูลเชิงจำแนก และภาพรวมการทดสอบไคกำลังสอง (Chi-Square Test)}

\textbf{ข้อมูลเชิงจำแนก (Categorical data)} คือข้อมูลเชิงคุณภาพที่แบ่งเป็นมาตรานามบัญญัติ (Nominal Scale เช่น เพศ, สาขาที่เรียน) หรือมาตราเรียงลำดับ (Ordinal Scale เช่น ระดับความพึงพอใจ 1--5) การทดสอบไคกำลังสองใช้กับข้อมูลลักษณะนี้เป็นหลัก โดยแบ่งเป็น 3 ประเภทตามลักษณะคำถาม
\begin{center}\small
\begin{tabular}{@{}p{3.3cm}p{9.8cm}@{}}
\toprule
\rowcolor{lightbg}\textbf{ประเภทการทดสอบ} & \textbf{ใช้เมื่อ}\\
\midrule
การทดสอบความเป็นอิสระกัน (Test of Independence) & มีตัวแปรเชิงคุณภาพ 2 ตัว วัดจาก\textbf{ประชากรกลุ่มเดียว} ต้องการทราบว่าตัวแปรทั้งสองสัมพันธ์กันหรือไม่\\
การทดสอบความคล้ายคลึงกัน (Test of Homogeneity) & มี\textbf{ตั้งแต่ 2 ประชากรอิสระกัน}ขึ้นไป แต่ละประชากรจำแนกด้วยตัวแปรเชิงคุณภาพ 1 ตัว ต้องการทราบว่าสัดส่วนของแต่ละคุณลักษณะในทุกประชากรเท่ากันหรือไม่\\
การทดสอบเทียบความกลมกลืนกัน (Goodness-of-Fit Test) & มีตัวอย่างกลุ่มเดียว จำแนกเป็นหลายคุณลักษณะ ต้องการทราบว่าการแจกแจงของข้อมูลตรงกับที่คาดไว้ (เช่น การแจกแจงทฤษฎีหนึ่ง ๆ) หรือไม่\\
\bottomrule
\end{tabular}
\end{center}
\concept{ทั้ง 3 แบบใช้สูตรตัวสถิติทดสอบ $\chi^2$ เหมือนกันทุกประการ (เปรียบเทียบความถี่ที่สังเกตได้ $O$ กับความถี่ที่คาดหวัง $E$) \textbf{ความแตกต่างอยู่ที่การออกแบบเก็บข้อมูลและคำถามวิจัย} ไม่ใช่สูตรคำนวณ: ถ้าสุ่มตัวอย่างมา\textbf{1 กลุ่ม}แล้ววัด 2 ตัวแปรพร้อมกัน $\to$ ถามว่า ``สัมพันธ์กันไหม'' (Independence); ถ้าสุ่มตัวอย่างมาจาก\textbf{หลายประชากรที่กำหนดขนาดกลุ่มไว้ล่วงหน้า} แล้ววัด 1 ตัวแปร $\to$ ถามว่า ``สัดส่วนเท่ากันไหม'' (Homogeneity); ถ้าสุ่มตัวอย่างมา 1 กลุ่มแล้วดูว่าตรงกับการแจกแจงที่คาดไว้หรือไม่ $\to$ (Goodness-of-Fit) นี่คือเหตุผลที่ตัวอย่างกรณีศึกษาเดียวกันสามารถนำไปใช้ตอบได้ทั้งคำถามแบบ Independence และแบบ Homogeneity โดยได้ค่าสถิติทดสอบชุดเดียวกัน (ดูตัวอย่างกรณีศึกษาด้านล่าง)}

\clearpage
\section{การทดสอบความเป็นอิสระกัน (Test of Independence)}

\textbf{แนวคิด:} ทดสอบว่าตัวแปรเชิงคุณภาพ 2 ตัวแปรที่วัดจากตัวอย่างกลุ่มเดียวกัน มีความสัมพันธ์กันหรือไม่ (เช่น เพศ กับ ระดับความเครียด) ถ้าผลทดสอบสรุปว่าเป็นอิสระต่อกัน แปลว่านำตัวแปรหนึ่งไปอธิบายอีกตัวแปรหนึ่งไม่ได้ แต่ถ้าไม่เป็นอิสระต่อกัน แปลว่าตัวแปรทั้งสองมีความสัมพันธ์กัน

\dist{ตารางการณ์จร (Contingency Table) ขนาด $r\times c$}
\begin{center}\small
\begin{tabular}{@{}c|cccc|c@{}}
\toprule
ตัวแปร A \textbackslash\ ตัวแปร B & $B_1$ & $B_2$ & $\dots$ & $B_c$ & รวม\\
\midrule
$A_1$ & $O_{11}$ & $O_{12}$ & $\dots$ & $O_{1c}$ & $O_{1.}$\\
$A_2$ & $O_{21}$ & $O_{22}$ & $\dots$ & $O_{2c}$ & $O_{2.}$\\
$\vdots$ & $\vdots$ & $\vdots$ & & $\vdots$ & $\vdots$\\
$A_r$ & $O_{r1}$ & $O_{r2}$ & $\dots$ & $O_{rc}$ & $O_{r.}$\\
\midrule
รวม & $O_{.1}$ & $O_{.2}$ & $\dots$ & $O_{.c}$ & $O_{..}=n$\\
\bottomrule
\end{tabular}
\end{center}

\begin{formulabox}
$$H_0:\text{ตัวแปร A และ B เป็นอิสระกัน}\qquad H_1:\text{ตัวแปร A และ B ไม่เป็นอิสระกัน}$$
ภายใต้ $H_0$ (เป็นอิสระกัน): $P(A_i\cap B_j)=P(A_i)P(B_j)$ ประมาณด้วยสัดส่วนตัวอย่าง $O_{i.}/n$ และ $O_{.j}/n$ จะได้ความถี่คาดหวัง
$$E_{ij}=nP(A_i\cap B_j)=n\left(\frac{O_{i.}}{n}\right)\left(\frac{O_{.j}}{n}\right)=\frac{O_{i.}\times O_{.j}}{O_{..}}$$
$$\chi^2=\sum_{i=1}^r\sum_{j=1}^c\frac{(O_{ij}-E_{ij})^2}{E_{ij}} \sim \chi^2_{(v)},\qquad v=(r-1)(c-1)$$
\end{formulabox}
\concept{สูตร $E_{ij}$ มาจากการสมมติว่า $H_0$ เป็นจริง (ตัวแปรเป็นอิสระกัน) แล้วประมาณความน่าจะเป็นของแต่ละแถว/คอลัมน์ด้วยสัดส่วนที่สังเกตได้จริง เพราะเราไม่ทราบ $P(A_i)$ และ $P(B_j)$ ที่แท้จริงของประชากร นี่คือหลักการเดียวกับการประมาณค่าพารามิเตอร์ด้วยสถิติจากตัวอย่าง (point estimation) ที่พบในบทก่อนหน้า}

\dist{กฎการยุบเซลล์ (Cell Merging Rule)}
ถ้ามีจำนวนเซลล์ที่ $E_{ij}<5$ \textbf{เกินร้อยละ 20} ของจำนวนเซลล์ทั้งหมด จำเป็นต้องยุบ (รวม) แถวหรือคอลัมน์ที่เกี่ยวข้องเข้าด้วยกัน เพื่อให้ค่าคาดหวังในแต่ละเซลล์มากพอที่การประมาณด้วยการแจกแจงไคสแควร์จะแม่นยำ
\concept{นี่เป็นเงื่อนไขเดียวกับที่ SPSS รายงานอัตโนมัติท้ายตาราง Chi-Square Tests เช่น ``N cells (\%) have expected count less than 5'' -- ถ้าเปอร์เซ็นต์ที่ SPSS รายงานเกิน 20\% ต้องกลับไปยุบเซลล์ก่อนแล้วจึงทดสอบใหม่ นอกจากนี้ ถ้าหลังยุบเซลล์แล้วเหลือตารางขนาด $2\times2$ พอดี และมีบางเซลล์ที่ $5\le E_{ij}\le10$ จะต้องใช้ \textbf{Yate's Correction for Continuity} (ลบ 0.5 ออกจากค่าสัมบูรณ์ของผลต่างก่อนยกกำลังสอง) เพื่อลดความเอนเอียงจากการประมาณค่าต่อเนื่องด้วยการแจกแจงไม่ต่อเนื่อง}

\clearpage
\section{ตัวอย่างที่ 1 (Independence): ภาวะการเป็นผู้นำ}

\begin{qbox}
ในการศึกษาภาวะการเป็นผู้นำ โดยสุ่มผู้บริหารขององค์กรแห่งหนึ่งมาจำนวน 75 คน แล้ววัดคะแนนภาวะการเป็นผู้นำ แบ่งเป็นระดับ ``มาก'' และ ``ปานกลางหรือน้อย'' อยากทราบว่า เพศ กับ ภาวะการเป็นผู้นำ มีความสัมพันธ์กันหรือไม่ ที่ระดับนัยสำคัญ 0.05
\end{qbox}

\begin{stepbox}{ขั้นที่ 1--2: สมมติฐาน และระดับนัยสำคัญ}
$$H_0:\text{เพศและภาวะการเป็นผู้นำไม่มีความสัมพันธ์กัน}\qquad H_1:\text{เพศและภาวะการเป็นผู้นำมีความสัมพันธ์กัน}$$
กำหนดระดับนัยสำคัญ $\alpha=0.05$
\concept{ข้อมูลนี้มาจากตัวอย่าง\textbf{กลุ่มเดียว} (ผู้บริหาร 75 คน) ที่วัด 2 ตัวแปรพร้อมกัน (เพศ และ ภาวะผู้นำ) เข้าเงื่อนไขของการทดสอบความเป็นอิสระกันพอดี}
\end{stepbox}

\begin{stepbox}{ขั้นที่ 3: คำนวณความถี่คาดหวัง}
ข้อมูลดิบแบ่งภาวะผู้นำเป็น 3 ระดับ (มาก/ปานกลาง/น้อย) โดยคอลัมน์รวมเท่ากับ 22, 46 และ 7 ตามลำดับ เมื่อคำนวณ $E_{ij}$ พบว่า $E_{13}=4.1$ และ $E_{23}=2.9$ ซึ่งน้อยกว่า 5 ทั้งคู่ (2 ใน 6 เซลล์ $=33.3\%>20\%$) จึงต้อง\textbf{ยุบคอลัมน์ ``ปานกลาง'' กับ ``น้อย'' เข้าด้วยกัน}เป็น ``ปานกลางหรือน้อย'' ได้ตารางใหม่ขนาด $2\times2$ ดังนี้

\begin{center}\small
\begin{tabular}{@{}lccc@{}}
\toprule
เพศ & มาก & ปานกลางหรือน้อย & รวม\\
\midrule
ชาย & 15 & 29 & 44\\
หญิง & 7 & 24 & 31\\
\midrule
รวม & 22 & 53 & 75\\
\bottomrule
\end{tabular}
\end{center}

$$E_{11}=\frac{44\times22}{75}=12.91 \quad E_{12}=\frac{44\times53}{75}=31.09 \quad E_{21}=\frac{31\times22}{75}=9.09 \quad E_{22}=\frac{31\times53}{75}=21.91$$
\concept{หลังยุบเป็น $2\times2$ แล้ว ทุกเซลล์มี $E_{ij}\ge5$ (น้อยสุดคือ 9.09) เข้าเงื่อนไข 20\% เรียบร้อย แต่เนื่องจากบางเซลล์อยู่ในช่วง $5\le E_{ij}\le10$ (คือ $E_{11}=12.91$ ไม่เข้าเงื่อนไข แต่ $E_{21}=9.09$ เข้าเงื่อนไขนี้) จึงต้องใช้ Yate's Correction for Continuity ในการคำนวณ $\chi^2$ ขั้นถัดไป}
\end{stepbox}

\begin{stepbox}{ขั้นที่ 4: คำนวณค่าสถิติทดสอบ (พร้อม Yate's Correction)}
$$\chi^2=\sum_{i=1}^2\sum_{j=1}^2\frac{\left(|O_{ij}-E_{ij}|-0.5\right)^2}{E_{ij}}$$
$$\chi^2=\frac{(|15-12.91|-0.5)^2}{12.91}+\frac{(|29-31.09|-0.5)^2}{31.09}+\frac{(|7-9.09|-0.5)^2}{9.09}+\frac{(|24-21.91|-0.5)^2}{21.91}=0.67$$
$$v=(r-1)(c-1)=(2-1)(2-1)=1 \qquad p\text{-value}=P(\chi^2_{(1)}>0.67)=0.413$$
\end{stepbox}

\begin{stepbox}{ขั้นที่ 5: ทำการตัดสินใจ และสรุปผลการทดสอบ}
p-value $=0.413>\alpha=0.05\Rightarrow$ \textbf{ยอมรับ} $H_0$

สรุป: เพศและภาวะการเป็นผู้นำ \textbf{ไม่มีความสัมพันธ์กัน} ที่ระดับนัยสำคัญ 0.05
\end{stepbox}

\clearpage
\section{ตัวอย่างกรณีศึกษา (Independence): เพศ กับ ระดับความเครียด}

\begin{qbox}
จากกรณีศึกษาความเครียดของนักศึกษาจากโควิด-19 ($n=150$) โดยแบ่งระดับความเครียดเป็น 3 ระดับจากคะแนน: 0--20 = ต่ำ, 21--31 = ปานกลาง, 32--40 = สูง อยากทราบว่า เพศของนักศึกษา มีความสัมพันธ์กับระดับความเครียด หรือไม่ ที่ระดับนัยสำคัญ 0.05
\end{qbox}

\begin{stepbox}{ขั้นที่ 1--2: สมมติฐาน และระดับนัยสำคัญ}
$$H_0:\text{เพศและระดับความเครียดไม่มีความสัมพันธ์กัน}\qquad H_1:\text{เพศและระดับความเครียดมีความสัมพันธ์กัน}$$
กำหนดระดับนัยสำคัญ $\alpha=0.05$; คำสั่งใน SPSS: \texttt{Analyze > Descriptive Statistics > Crosstabs...} (ใส่เพศและระดับความเครียดในช่อง Row/Column แล้วเลือก Statistics $\to$ Chi-square, Cells $\to$ Expected)
\end{stepbox}

\begin{stepbox}{ขั้นที่ 3--4: ตารางการณ์จร และคำนวณค่าสถิติทดสอบ}
\spssout{เพศ * ระดับความเครียด Crosstabulation}
\begin{center}\small
\begin{tabular}{@{}llcccc@{}}
\toprule
& & ระดับต่ำ & ระดับปานกลาง & ระดับสูง & Total\\
\midrule
\multirow{2}{*}{ชาย} & Count & 35 & 59 & 12 & 106\\
& Expected Count & 38.9 & 56.5 & 10.6 & 106.0\\
\multirow{2}{*}{หญิง} & Count & 20 & 21 & 3 & 44\\
& Expected Count & 16.1 & 23.5 & 4.4 & 44.0\\
\midrule
\multirow{2}{*}{Total} & Count & 55 & 80 & 15 & 150\\
& Expected Count & 55.0 & 80.0 & 15.0 & 150.0\\
\bottomrule
\end{tabular}
\end{center}

\spssout{Chi-Square Tests}
\begin{center}\small
\begin{tabular}{@{}lccc@{}}
\toprule
& Value & df & Asymp. Sig. (2-sided)\\
\midrule
Pearson Chi-Square & $2.309^a$ & 2 & .315\\
Likelihood Ratio & 2.313 & 2 & .315\\
N of Valid Cases & 150 & & \\
\bottomrule
\end{tabular}
\end{center}
{\scriptsize $a$. 1 cell (16.7\%) have expected count less than 5. The minimum expected count is 4.40.}
\concept{หมายเหตุท้ายตาราง ``1 cell (16.7\%)'' คือจุดที่ต้องตรวจก่อนอ่านค่า Sig. เสมอ: มีเพียง 1 ใน 6 เซลล์ $=16.7\%<20\%$ จึง\textbf{ไม่ต้องยุบเซลล์}และอ่านค่า Pearson Chi-Square ได้โดยตรง ต่างจากตัวอย่างภาวะผู้นำก่อนหน้าที่เกิน 20\% จึงต้องยุบเซลล์ก่อน}
\end{stepbox}

\begin{stepbox}{ขั้นที่ 5: ทำการตัดสินใจ และสรุปผลการทดสอบ}
p-value $=.315>\alpha=0.05\Rightarrow$ \textbf{ยอมรับ} $H_0$

สรุป: เพศและระดับความเครียด \textbf{ไม่มีความสัมพันธ์กัน} ($p>0.05$)
\end{stepbox}

\clearpage
\section{การทดสอบความคล้ายคลึงกัน (Test of Homogeneity)}

\textbf{แนวคิด:} ใช้เมื่อมี\textbf{ตั้งแต่ 2 ประชากรที่เป็นอิสระกัน}ขึ้นไป โดยแต่ละประชากรถูกจำแนกด้วยตัวแปรเชิงคุณภาพ 1 ตัว ออกเป็นอย่างน้อย 2 คุณลักษณะ ต้องการทดสอบว่าสัดส่วนของแต่ละคุณลักษณะใน\textbf{ทุกประชากรเท่ากันหรือไม่} ถ้าสัดส่วนไม่แตกต่างกัน แสดงว่าประชากรเหล่านั้น ``คล้ายคลึงกัน''

\begin{formulabox}
$$H_0:\text{สัดส่วนของแต่ละคุณลักษณะในทุกประชากรเท่ากัน (ประชากรคล้ายคลึงกัน)}$$
$$H_1:\text{มีอย่างน้อย 1 ประชากรที่สัดส่วนแตกต่างจากกลุ่มอื่น}$$
ใช้ตัวสถิติทดสอบ $\chi^2$ สูตรเดียวกับ Test of Independence ทุกประการ
\end{formulabox}
\concept{\textbf{ความแตกต่างจาก Test of Independence อยู่ที่การออกแบบเก็บข้อมูล ไม่ใช่สูตรคำนวณ:} Independence สุ่มตัวอย่าง 1 กลุ่มมาวัด 2 ตัวแปรพร้อมกัน (ผลรวมแต่ละแถว/คอลัมน์เกิดขึ้นเองจากการสุ่ม) ส่วน Homogeneity \textbf{กำหนดขนาดตัวอย่างของแต่ละประชากรไว้ล่วงหน้า} (เช่น สุ่มสินค้าผลัดเช้า 950 ชิ้น ผลัดบ่าย 945 ชิ้น ผลัดดึก 940 ชิ้น -- ตัวเลขเหล่านี้ผู้วิจัยเป็นผู้กำหนด ไม่ได้เกิดจากการสุ่มเอง) แล้ววัด 1 ตัวแปรในแต่ละกลุ่ม คำถามวิจัยจึงเปลี่ยนจาก ``สัมพันธ์กันไหม'' เป็น ``สัดส่วนเท่ากันไหม'' แต่ในทางคณิตศาสตร์ทั้งสองแบบคำนวณ $E_{ij}$ และ $\chi^2$ ด้วยสูตรเดียวกันทุกประการ}

\clearpage
\section{ตัวอย่างที่ 2 (Homogeneity): คุณภาพสินค้าจำแนกตามผลัดการผลิต}

\begin{qbox}
ในการศึกษาคุณภาพของสินค้าที่ผลิตจากผลัดต่าง ๆ ของโรงงานแห่งหนึ่ง โดยสุ่มสินค้าที่ผลิตจากผลัดเช้า บ่าย และดึก มาจำนวน 950, 945 และ 940 ชิ้น ตามลำดับ จำแนกตามคุณภาพได้ดังนี้ จงทดสอบว่าคุณภาพสินค้าทั้ง 3 ผลัดมีสัดส่วนแตกต่างกันหรือไม่ ที่ระดับนัยสำคัญ 0.05

\begin{center}\small
\begin{tabular}{@{}lccc@{}}
\toprule
ผลัดการผลิต & ชำรุด & ดี & รวม\\
\midrule
ผลัดเช้า & 45 & 905 & 950\\
ผลัดบ่าย & 55 & 890 & 945\\
ผลัดดึก & 70 & 870 & 940\\
\midrule
รวม & 170 & 2{,}665 & 2{,}835\\
\bottomrule
\end{tabular}
\end{center}
\end{qbox}

\begin{stepbox}{ขั้นที่ 1--3: สมมติฐาน ระดับนัยสำคัญ และความถี่คาดหวัง}
$$H_0:\text{สัดส่วนสินค้าชำรุดของทั้ง 3 ผลัดเท่ากัน}\qquad H_1:\text{มีอย่างน้อย 1 ผลัดที่สัดส่วนแตกต่างจากผลัดอื่น}$$
กำหนดระดับนัยสำคัญ $\alpha=0.05$
$$E_{ij}=\frac{O_{i.}\times O_{.j}}{O_{..}}$$
$$E_{11}=\frac{950\times170}{2{,}835}=56.97 \quad E_{21}=\frac{945\times170}{2{,}835}=56.67 \quad E_{31}=\frac{940\times170}{2{,}835}=56.37$$
$$E_{12}=\frac{950\times2{,}665}{2{,}835}=893.03 \quad E_{22}=\frac{945\times2{,}665}{2{,}835}=888.33 \quad E_{32}=\frac{940\times2{,}665}{2{,}835}=883.63$$
\concept{ทุกเซลล์มี $E_{ij}\gg5$ (ค่าต่ำสุด $=56.37$) จึง\textbf{ไม่ต้องยุบเซลล์} สังเกตว่าถึงแม้จำนวนชิ้นชำรุดที่สุ่มมา (45, 55, 70) จะดูน้อย แต่เพราะขนาดตัวอย่างรวมของแต่ละผลัดใหญ่มาก ($\approx950$ ชิ้น) ความถี่คาดหวังจึงยังคงมากกว่า 5 อยู่มาก}
\end{stepbox}

\begin{stepbox}{ขั้นที่ 4: คำนวณค่าสถิติทดสอบ}
$$v=(r-1)(c-1)=(3-1)(2-1)=2$$
$$\chi^2=\sum_{i=1}^3\sum_{j=1}^2\frac{(O_{ij}-E_{ij})^2}{E_{ij}}=\frac{(45-56.97)^2}{56.97}+\cdots+\frac{(870-883.63)^2}{883.63}=6.23$$

\excelout{หาค่า p-value ด้วย \texttt{=CHIDIST(6.23,2)}}
$$p\text{-value}=P(\chi^2_{(2)}>6.23)=0.0443$$
\end{stepbox}

\begin{stepbox}{ขั้นที่ 5: ทำการตัดสินใจ และสรุปผลการทดสอบ}
p-value $=0.0443<\alpha=0.05\Rightarrow$ \textbf{ปฏิเสธ} $H_0$

สรุป: คุณภาพของสินค้าที่ผลิตจากทั้ง 3 ผลัด มีสัดส่วนที่\textbf{แตกต่างกัน} ที่ระดับนัยสำคัญ 0.05
\concept{เปรียบเทียบสัดส่วนสินค้าชำรุดคร่าว ๆ: เช้า $45/950=4.7\%$, บ่าย $55/945=5.8\%$, ดึก $70/940=7.4\%$ -- แนวโน้มคือกะดึกมีสัดส่วนสินค้าชำรุดสูงกว่ากะอื่นตามลำดับ ซึ่งสอดคล้องกับสามัญสำนึกว่าการทำงานกะดึกอาจมีผลต่อคุณภาพการผลิต และผล $\chi^2$ ที่มีนัยสำคัญยืนยันว่าความแตกต่างนี้ไม่ได้เกิดจากความบังเอิญในการสุ่มตัวอย่างเพียงอย่างเดียว}
\end{stepbox}

\clearpage
\section{ตัวอย่างกรณีศึกษา (Homogeneity): เพศ กับ ระดับความเครียด}

\begin{qbox}
จากกรณีศึกษาความเครียดของนักศึกษาจากโควิด-19 ชุดเดียวกับตัวอย่าง Independence ก่อนหน้า อยากทราบว่า นักศึกษาชาย และนักศึกษาหญิง มีสัดส่วนของระดับความเครียด แตกต่างกันหรือไม่ ที่ระดับนัยสำคัญ 0.05
\end{qbox}

เนื่องจากใช้ข้อมูลชุดเดียวกัน (เพศ $\times$ ระดับความเครียด, $n=150$) ตัวสถิติทดสอบที่คำนวณได้จึง\textbf{เป็นชุดเดียวกัน}กับตัวอย่างกรณีศึกษา Independence ก่อนหน้าทุกประการ: Pearson Chi-Square $=2.309$, $df=2$, p-value $=.315$

p-value $=.315>\alpha=0.05\Rightarrow$ \textbf{ยอมรับ} $H_0$

สรุป: นักศึกษาชาย และนักศึกษาหญิง มีสัดส่วนของระดับความเครียด\textbf{ไม่แตกต่างกัน} ($p>0.05$)
\concept{ตัวอย่างนี้จงใจใช้\textbf{ข้อมูลชุดเดียวกัน}กับตัวอย่าง Independence เพื่อสาธิตประเด็นสำคัญที่สุดของบทนี้: คำถามวิจัย ``เพศสัมพันธ์กับความเครียดไหม'' (Independence) กับคำถาม ``สัดส่วนความเครียดของชายกับหญิงต่างกันไหม'' (Homogeneity) ฟังดูต่างกัน แต่เมื่อข้อมูลมาจากตารางการณ์จรเดียวกัน ตัวสถิติทดสอบและข้อสรุปทางคณิตศาสตร์จะ\textbf{เหมือนกันเป๊ะ} -- สิ่งที่ต่างกันจริง ๆ คือ \textit{มุมมองการตั้งคำถามและบริบทเชิงเนื้อหา}เท่านั้น ในทางปฏิบัติ ถ้าการเก็บข้อมูลไม่ได้กำหนดขนาดกลุ่มตัวอย่างของแต่ละเพศไว้ล่วงหน้า (สุ่มมาจากประชากรเดียวแล้วเพศเกิดขึ้นเอง) จะเรียกว่าเป็นการทดสอบความเป็นอิสระกันอย่างเคร่งครัดกว่า}

\clearpage
\section{การทดสอบเทียบความกลมกลืนกัน (Goodness-of-Fit Test)}

\textbf{แนวคิด:} ใช้ทดสอบว่าการแจกแจงของข้อมูลตัวอย่างกลุ่มเดียว (ที่จำแนกได้ตั้งแต่ 2 คุณลักษณะขึ้นไป) สอดคล้องกับการแจกแจงที่คาดหวังไว้ (เช่น การแจกแจงทวินาม, ปัวซง, หรือสัดส่วนทางทฤษฎีใด ๆ) หรือไม่

\begin{formulabox}
$$H_0:\text{ข้อมูลมีการแจกแจงตามที่กำหนด}\qquad H_1:\text{ข้อมูลไม่มีการแจกแจงตามที่กำหนด}$$
$$E_i=np_i \qquad \chi^2=\sum_{i=1}^{k}\frac{(O_i-E_i)^2}{E_i}\sim\chi^2_{(v)}$$
$$v=k-1-m \quad (m=\text{จำนวนพารามิเตอร์ที่ต้องประมาณจากข้อมูลตัวอย่าง})$$
\end{formulabox}
\concept{ต่างจาก Independence/Homogeneity ตรงที่ Goodness-of-Fit เปรียบเทียบกับ\textbf{การแจกแจงทางทฤษฎี} ไม่ใช่เปรียบเทียบระหว่างกลุ่มตัวอย่างด้วยกันเอง และ $df$ ต้อง\textbf{หักลบจำนวนพารามิเตอร์ที่ประมาณจากข้อมูล} ($m$) เพิ่มเติมจากปกติ ($k-1$) เช่น ถ้าทดสอบการแจกแจงปัวซงและต้องประมาณค่าเฉลี่ย $\lambda$ จากข้อมูลตัวอย่างเอง จะต้องหัก $df$ ออกอีก 1 หน่วย}

\clearpage
\section{ตัวอย่าง (Goodness-of-Fit): จำนวนครั้งที่ร้องขอค่ารักษาพยาบาล}

\begin{qbox}
บริษัทประกันชีวิตแห่งหนึ่งสุ่มผู้ซื้อประกันที่มีอายุต่ำกว่า 50 ปี และซื้อประกันมาแล้วอย่างน้อย 4 ปี มาจำนวน 200 คน บันทึกจำนวนครั้งที่ร้องขอค่ารักษาพยาบาลในช่วง 4 ปีที่ผ่านมา ได้ผลดังนี้ ที่ระดับนัยสำคัญ 0.05 บริษัทจะสรุปได้หรือไม่ว่าจำนวนครั้งที่ร้องขอมีการแจกแจงแบบปัวซง (Poisson Distribution)

\begin{center}\small
\begin{tabular}{@{}lccccccccc@{}}
\toprule
จำนวนครั้งที่ร้องขอ ($x$) & 0 & 1 & 2 & 3 & 4 & 5 & 6 & 7\\
\midrule
จำนวนผู้ซื้อประกัน ($O_i$) & 22 & 53 & 58 & 39 & 20 & 5 & 2 & 1\\
\bottomrule
\end{tabular}
\end{center}
\end{qbox}

\begin{stepbox}{ขั้นที่ 1--3: สมมติฐาน ระดับนัยสำคัญ และประมาณค่าพารามิเตอร์}
$$H_0:\text{จำนวนครั้งที่ร้องขอมีการแจกแจงแบบปัวซง}\qquad H_1:\text{ไม่มีการแจกแจงแบบปัวซง}$$
กำหนดระดับนัยสำคัญ $\alpha=0.05$

เนื่องจากไม่ทราบค่าพารามิเตอร์ $\lambda$ ของการแจกแจงปัวซง $f(x;\lambda)=\dfrac{e^{-\lambda}\lambda^x}{x!}$ (ซึ่งมี $E(X)=\lambda$) จึงประมาณด้วยค่าเฉลี่ยตัวอย่าง
$$\bar X=\frac{\sum fX}{\sum f}=\frac{(0)(22)+(1)(53)+\cdots+(7)(1)}{200}=2.05$$
\concept{นี่คือจุดสำคัญที่ทำให้ตัวอย่างนี้ต่าง จากตัวอย่างการทดสอบ Independence/Homogeneity ก่อนหน้า: การแจกแจงปัวซงมีพารามิเตอร์ $\lambda$ ที่ต้อง\textbf{ประมาณจากข้อมูลตัวอย่างเอง} ($\hat\lambda=\bar X=2.05$) ทำให้ต้องหัก $df$ เพิ่มอีก 1 หน่วยตอนคำนวณ p-value ในขั้นถัดไป}
\end{stepbox}

\begin{stepbox}{ขั้นที่ 4: คำนวณความถี่คาดหวัง ยุบเซลล์ และหาค่าสถิติทดสอบ}
ภายใต้ $H_0$ คำนวณ $p_i=\dfrac{e^{-2.05}(2.05)^x}{x!}$ และ $E_i=np_i=200p_i$ สำหรับ $x=0,1,\dots,7$

\begin{center}\small
\begin{tabular}{@{}lcccccccc@{}}
\toprule
$x$ & 0 & 1 & 2 & 3 & 4 & 5 & 6 & $\ge7$\\
\midrule
$O_i$ & 22 & 53 & 58 & 39 & 20 & 5 & 2 & 1\\
$p_i$ & .1287 & .2639 & .2705 & .1848 & .0947 & .0388 & .0133 & .0053\\
$E_i$ & 25.74 & 52.78 & 54.10 & 36.96 & 18.94 & 7.76 & 2.66 & 1.06\\
\bottomrule
\end{tabular}
\end{center}

พบว่าเซลล์ $x=6$ และ $x\ge7$ มี $E_i<5$ ทั้งคู่ (คิดเป็น $2/8=25\%>20\%$ ของจำนวนเซลล์ทั้งหมด) จึงต้อง\textbf{ยุบรวม} $x=6$ กับ $x\ge7$ เข้าด้วยกันเป็นหมวด ``$x\ge6$'': $O=2+1=3$, $E=2.66+1.06=3.72$ ทำให้เหลือทั้งหมด 7 หมวด

$$v=k-1-m=7-1-1=5 \quad(\text{หัก 1 เพราะประมาณ }\lambda\text{ จากข้อมูล})$$
$$\chi^2=\sum_{i}\frac{(O_i-E_i)^2}{E_i}=\frac{(22-25.74)^2}{25.74}+\cdots+\frac{(3-3.72)^2}{3.72}=2.118$$

\excelout{หาค่า p-value ด้วย \texttt{=CHIDIST(2.118,5)}}
$$p\text{-value}=P(\chi^2_{(5)}>2.118)=0.8326$$
\end{stepbox}

\begin{stepbox}{ขั้นที่ 5: ทำการตัดสินใจ และสรุปผลการทดสอบ}
p-value $=0.8326>\alpha=0.05\Rightarrow$ \textbf{ยอมรับ} $H_0$

สรุป: จำนวนครั้งที่ผู้ซื้อประกันร้องขอค่ารักษาพยาบาลในช่วง 4 ปีที่ผ่านมา \textbf{มีการแจกแจงแบบปัวซง} ที่ระดับนัยสำคัญ 0.05
\concept{p-value ที่สูงมาก ($0.83$) แปลว่าข้อมูลที่สังเกตได้ ($O_i$) ใกล้เคียงกับค่าที่คาดหวังภายใต้การแจกแจงปัวซง ($E_i$) มาก จนแทบไม่มีหลักฐานคัดค้าน $H_0$ เลย ต่างจากตัวอย่าง Homogeneity ก่อนหน้าที่ p-value ใกล้ $\alpha$ ($0.044$) ซึ่งต้องระมัดระวังในการตีความมากกว่า}
\end{stepbox}

\clearpage
\section{สรุปการเลือกใช้ ANOVA และ Chi-square Test}

\begin{center}\small
\begin{tabular}{@{}p{5.3cm}p{7.6cm}@{}}
\toprule
\textbf{สถานการณ์} & \textbf{วิธีทดสอบ}\\\midrule
เปรียบเทียบค่าเฉลี่ยของ $\ge3$ ประชากรอิสระกัน & One-Way ANOVA ($F$-test) ตามด้วย Multiple Comparison (LSD/Tukey/Scheffe/Bonferroni) ถ้าปฏิเสธ $H_0$\\
ตัวแปรเชิงคุณภาพ 2 ตัว วัดจากตัวอย่าง 1 กลุ่ม -- สัมพันธ์กันไหม & Chi-square Test of Independence\\
ตัวแปรเชิงคุณภาพ 1 ตัว วัดจากหลายประชากรอิสระ (กำหนดขนาดกลุ่มไว้ก่อน) -- สัดส่วนเท่ากันไหม & Chi-square Test of Homogeneity\\
ตัวอย่าง 1 กลุ่ม จำแนกหลายคุณลักษณะ -- ตรงกับการแจกแจงที่คาดไว้ไหม & Chi-square Goodness-of-Fit Test\\
\bottomrule
\end{tabular}
\end{center}

\begin{formulabox}
\textbf{หลักทั่วไปที่ใช้ซ้ำทุกครั้ง (5 ขั้นตอน เหมือนบท Part 1 ทุกประการ):}
\begin{enumerate}
\item ตั้ง $H_0,H_1$ -- ดูว่าโจทย์ถามเรื่องค่าเฉลี่ย ($\ge3$ กลุ่ม $\to$ ANOVA) หรือความถี่/สัดส่วน/ความสัมพันธ์ ($\to$ Chi-square และเลือกประเภทให้ถูกจากการออกแบบเก็บข้อมูล)
\item กำหนด $\alpha$ -- โดยทั่วไป 0.05
\item ตรวจข้อตกลงเบื้องต้น/เตรียมข้อมูล -- ANOVA ต้องตรวจการแจกแจงปกติและความแปรปรวนเท่ากันก่อน; Chi-square ต้องตรวจกฎการยุบเซลล์ ($E_{ij}<5$ เกิน 20\%) ก่อนเสมอ
\item คำนวณค่าสถิติ/p-value -- อ่านจาก Output SPSS/Excel ให้ถูกตาราง ถูกแถว (ANOVA: ปฏิเสธแล้วต้องทำ Multiple Comparison ต่อ)
\item เทียบ p-value กับ $\alpha$ แล้วสรุปผลกลับไปที่คำถามของโจทย์ (ANOVA บอกว่าคู่ไหนต่างกันด้วยเสมอ; Chi-square บอกทิศทาง/ขนาดของความสัมพันธ์หรือความแตกต่างของสัดส่วนประกอบ)
\end{enumerate}
\end{formulabox}

\end{document}

